\documentclass[12pt]{article}
\usepackage[utf8]{inputenc}
\usepackage[T1]{fontenc}
\usepackage{amsmath,amssymb}
\usepackage{geometry}
\geometry{margin=2.5cm}

\begin{document}

\noindent\underline{\textbf{T.\ fundamentală a aritmeticii}}

\bigskip
\noindent $(\forall)\,a\in\mathbb{N}$, $a>1$, $\Rightarrow a=p_1\cdots p_n$

\noindent [se reprezintă în mod unic ca produs de nr.\ prime]

\bigskip
\noindent\underline{Lemă}: $(\forall)\,a\in\mathbb{N}$, $a>1$ \ $\exists\,p$=prim, $p\mid a$.

\bigskip
\noindent $P=\{d \mid d\mid a \text{ şi } d\neq 1\}$

\bigskip
\noindent Fie $p\in P$ cel mai mic nr.\ din $P$.

\noindent\underline{$p$=prim}.

\bigskip
\noindent P.A.\ $p\neq$prim $\ \exists\,d_1$ a.î.\ $d_1\mid p$, $d_1\neq 1$ şi $d_1\neq p$.

\[
\left.\begin{aligned} d_1&\mid p\\ p&\mid a\end{aligned}\right\} \Rightarrow d_1\mid a
\]
\[
\left.\begin{aligned} &\Rightarrow d_1\in P\\ &d_1<p\end{aligned}\right\} \Rightarrow p\text{=prim.} \quad \times
\]

\bigskip
\noindent\underline{Dem.\ teor}:

\noindent Pp.\ nu e adev.

\noindent Notez $P=\big\{$mulţ.\ nr.\ nat.\ (mai mari ca 1)\footnote{abrevierea din original (,,not.\ core'') e neclară -- reconstituită aici ca o condiție ,,$>1$'', potrivit cu enunțul teoremei.} ce nu se reprezintă ca

\noindent produs finit de nr.\ prime$\big\} \neq \varnothing$.

\bigskip
\noindent Fie $a$ = cel mai mic nr.\ din $P$.

\noindent $\Rightarrow a$ nu e prim

\bigskip
\noindent Fie $a=bc$.

\noindent $b,c\notin P$

\[
\Rightarrow b=p_1\cdots p_j
\]
\[
c=p_1'\cdots p_i'
\]
\[
\Rightarrow a = p_1\cdots p_j\cdot p_1'\cdots p_i'.
\]

\bigskip
\noindent\underline{Unicitatea}: \ $a = p_1\cdots p_n = p_1'\cdots p_m'$.

\[
\overset{(5)}{\Longrightarrow} p_1\mid p_1'\cdots p_m' \Rightarrow \exists\, i \text{ a.î. } p_1\mid p_i' \Rightarrow p_1=p_i'
\]

\end{document}
